\begin{proof}
Set
\[
  s_t:=t^{-(p+2)}M_p(t).
\]
Since the coefficient functions are bounded,
\(\|A_t\|_\infty=O(t^{-2})\) and
\(\|B_t\|_\infty=O(t^{-p}M_p(t))=o(t^{-2})\) by
\eqref{eq:moving-scale-growth}.  Also
\(\|e_t\|_\infty=o(t^{-2})\) by
\eqref{eq:moving-scale-error} and \eqref{eq:moving-scale-growth}.  Hence all
arguments of \(\Phi\) below lie in a fixed interval \((-1/2,1/2)\) for large
\(t\).

On this interval,
\[
  \Phi'(x)=\log(1+x),\qquad \Phi''(x)=\frac1{1+x},
\]
and \(|\Phi'(x)|\le C|x|\), \(|\Phi''(x)|\le C\).  The error \(e_t\) is
negligible because the mean-value theorem gives
\[
\begin{aligned}
\left|\int_E\{\Phi(A_t+B_t+e_t)-\Phi(A_t+B_t)\}\,\dd\lambda\right|
&\le
C(\|A_t\|_\infty+\|B_t\|_\infty+\|e_t\|_\infty)\|e_t\|_{L^1(\lambda)}\\
&=o(s_t),
\end{aligned}
\]
using \(\|A_t\|_\infty=O(t^{-2})\),
\(\|B_t\|_\infty+\|e_t\|_\infty=o(t^{-2})\), and
\(\|e_t\|_{L^1}=o(s_t/t^{-2})\).

It remains to compare \(A_t+B_t\) with \(A_t\).  Taylor's formula at \(A_t\)
gives
\[
\Phi(A_t+B_t)-\Phi(A_t)
=\Phi'(A_t)B_t+O(B_t^2)
\]
uniformly.  Since \(b\) is bounded,
\[
  \int_E O(B_t^2)\,\dd\lambda
  =O(t^{-2p}M_p(t)^2)
  =o(s_t),
\]
again by \(M_p(t)=o(t^{p-2})\).  Therefore
\[
  \int_E\{\Phi(A_t+B_t)-\Phi(A_t)\}\,\dd\lambda
  =
  t^{-p}M_p(t)\int_E\log(1+A_t)b\,\dd\lambda+o(s_t).
\]
Finally \(\log(1+A_t)=A_t+O(A_t^2)\) uniformly, so
\[
\begin{aligned}
t^{-p}M_p(t)\int_E\log(1+A_t)b\,\dd\lambda
&=
t^{-p}M_p(t)\int_EA_tb\,\dd\lambda+O(t^{-p}M_p(t)\|A_t\|_\infty^2)\\
&=
t^{-(p+2)}M_p(t)\int_Ea_2b\,\dd\lambda+o(s_t).
\end{aligned}
\]
The terms \(a_jt^{-j}\) with \(j\ge3\) contribute \(O(t^{-(p+3)}M_p(t))\),
and the quadratic logarithm remainder contributes \(O(t^{-(p+4)}M_p(t))\).
Combining the three estimates proves \eqref{eq:moving-scale-transfer}.
\end{proof}

\begin{proposition}[Finite density jet entropy transfer]%
\label{prop:finite-density-jet-entropy-transfer}
Let \(N\ge2\), put \(L=2N-2\), and let \((E,\mathcal E,\lambda)\) be a
probability space.  Suppose
\[
\delta_t=d_t+\rho_t,\qquad
d_t(y)=\sum_{n=2}^{L}a_n(y)t^{-n},
\]
where \(a_2,\ldots,a_L\in L^\infty(\lambda)\),
\(\|d_t\|_\infty=O(t^{-2})\), and
\[
\|\rho_t\|_\infty+\|\rho_t\|_{L^1(\lambda)}=o(t^{-L}).
\]
Then, with \(\Phi(s)=(1+s)\log(1+s)-s\),
\[
\int_E\Phi(\delta_t)\,\dd\lambda
=
\sum_{r=2}^{N}\frac{(-1)^r}{r(r-1)}
\int_E d_t(y)^r\,\dd\lambda
+o(t^{-(L+2)}).
\]
\end{proposition}

\begin{proof}
Since \(\|d_t\|_\infty=O(t^{-2})\) and
\(\|\rho_t\|_\infty=o(t^{-L})\), also
\(\|\delta_t\|_\infty=O(t^{-2})\).  Hence, uniformly in \(y\),
\[
\Phi(\delta_t)=
\sum_{r=2}^{N}\frac{(-1)^r}{r(r-1)}\delta_t^r
+O(|\delta_t|^{N+1}).
\]
The Taylor tail is \(O(t^{-2N-2})=o(t^{-(L+2)})\).  For every
\(2\le r\le N\), the binomial expansion and one \(L^1(\lambda)\) placement
of the remainder give
\[
\|\delta_t^r-d_t^r\|_{L^1(\lambda)}
\le
C_r\sum_{\ell=1}^{r}
\|d_t\|_\infty^{r-\ell}
\|\rho_t\|_\infty^{\ell-1}
\|\rho_t\|_{L^1(\lambda)}
=o(t^{-(L+2)}).
\]
Thus the Taylor polynomial may be evaluated on the finite density polynomial
\(d_t\) through order \(L+2\), which proves the claim.
\end{proof}

\medskip
\phantomsection\label{par:coefficient-evaluability-scope}
\noindent\textbf{Finite-moment coefficient level.}
The finite-moment level for the hierarchy is summarized in
Subsection~\ref{sec:entropy-theorem-map}; the algebraic obstruction
behind the leading threshold is proved below in
Proposition~\ref{prop:poisson-kl-moment-threshold-sharpness}.
\medskip

\begin{theorem}[Finite-moment Poisson entropy transfer and coefficient formula]%
\label{thm:poisson-kl-even-orders}
Let \(N\ge2\), put \(L=2N-2\), and let \(\nu\) be a probability law on
\(\RR\) with finite mean \(\bar{\gamma}\) and finite centred \(L\)th absolute
moment.  Let
\[
\begin{aligned}
X_c&=\gamma-\bar{\gamma},\\
\mu_n&=\int_{\RR} X_c^n\,\dd\nu(\gamma)\qquad(2\le n\le L),\\
h_t&=P_t*\nu,\qquad
g_t(x)&=P_t(x-\bar{\gamma}),
\end{aligned}
\]
then
\begin{equation}\label{eq:poisson-kl-even-orders}
D_{\mathrm{KL}}(h_t\|g_t)
=\sum_{m=2}^{N}A_{2m}(\nu)t^{-2m}+o(t^{-2N}).
\end{equation}
For \(4\le k\le2N\),
\begin{equation}\label{eq:coefficient-identification-transfer-expansion}
A_k(\nu)=
\sum_{r=2}^{\lfloor k/2\rfloor}
\frac{(-1)^r}{r(r-1)}
\sum_{\substack{n_1+\cdots+n_r=k\\2\le n_j\le2N-2}}
\mu_{n_1}\cdots\mu_{n_r}
\int_{\RR}u_{n_1}\cdots u_{n_r}\,\dd\omega .
\end{equation}
All retained odd \(A_k\) vanish.
\end{theorem}

\begin{proof}
Write
\[
d_t(y)=\sum_{n=2}^{L}\mu_nu_n(y)t^{-n}.
\]
Lemma~\ref{lem:exact-moment-density-truncation} gives
\(\delta_t=d_t+\rho_t\),
\(\|\rho_t\|_\infty+\|\rho_t\|_{L^1(\omega)}=o(t^{-L})\).  The Cayley modes
\(u_2,\ldots,u_L\) are bounded, hence \(\|d_t\|_\infty=O(t^{-2})\).  Since
\[
g_t(\bar{\gamma}+ty)t\,\dd y=\omega(\dd y),
\]
and \(\int\delta_t\,\dd\omega=0\),
\[
D_{\mathrm{KL}}(h_t\|g_t)
=\int_{\RR}\Phi(\delta_t(y))\,\omega(\dd y),
\qquad
\Phi(s)=(1+s)\log(1+s)-s .
\]
Applying
Proposition~\ref{prop:finite-density-jet-entropy-transfer} with
\((E,\lambda)=(\RR,\omega)\) gives
\[
D_{\mathrm{KL}}(h_t\|g_t)
=
\sum_{r=2}^{N}\frac{(-1)^r}{r(r-1)}
\int_{\RR}\left(\sum_{n=2}^{L}\mu_nu_n(y)t^{-n}\right)^r\omega(\dd y)
+o(t^{-2N}).
\]
Extracting the Laurent coefficient of \(t^{-k}\) gives
\eqref{eq:coefficient-identification-transfer-expansion}, and the odd
orders vanish by Theorem~\ref{thm:universal-entropy-coefficients}.

The constant-term evaluation uses a single Laurent convention throughout:
\[
Q_n(z)=2^{-n}(-\ii)^n
\sum_{k=1}^{n}\binom{n-1}{k-1}
\bigl(z^k+(-1)^nz^{-k}\bigr),
\qquad
J(n_1,\ldots,n_r)=[z^0]\prod_{j=1}^r Q_{n_j}(z).
\]
Equivalently, \(Q_n(z)=\sum_kq_{n,k}z^k\), \(q_{n,0}=0\), and
\[
q_{n,k}=2^{-n}(-\ii)^n
\operatorname{sgn}(k)^{\,n\bmod2}
\binom{n-1}{|k|-1}\qquad(0<|k|\le n).
\]
Thus the signs of the negative odd modes are already contained in
\(q_{n,k}\).  For example, with \(\sigma^2=\mu_2\),
\[
A_4=\frac{\sigma^4}{8},\qquad
A_6=\frac{\sigma^6+6\mu_3^2-8\sigma^2\mu_4}{64};
\]
  the sixth-order arithmetic is displayed in Lemma ``Sixth-order coefficient
  arithmetic.''  The auxiliary order-eight row values are summarized in
  Appendix ``Finite Laurent certificate through order
  eight''; the coefficient
identification itself is the main-text finite-moment theorem, not an appendix
input.
\end{proof}

\begin{lemma}[Positive-margin quotient remainder and entropy transfer]
\label{lem:positive-margin-quotient-remainder}
\label{rem:positive-margin-retained-order}
\label{prop:density-to-entropy-descent}
\label{lem:closed-cayley-mode-constant-envelopes}
Let \(L\ge2\), \(0<\eta\le1\), and let \(\nu\) have mean
\(\bar{\gamma}\) and finite centred \((L+\eta)\)th absolute moment.  With
\[
\delta_t(y):=
\frac{(P_t*\nu)(\bar{\gamma}+ty)}{P_t(ty)}-1
\]
there is a remainder \(\rho_{L,t}^{(\eta)}\) such that
\begin{equation}\label{eq:positive-margin-density-remainder}
\delta_t(y)=
\sum_{n=2}^{L}\mu_nu_n(y)t^{-n}+\rho_{L,t}^{(\eta)}(y),
\qquad
\|\rho_{L,t}^{(\eta)}\|_{\infty}
+\|\rho_{L,t}^{(\eta)}\|_{L^1(\omega)}
=O_\nu(t^{-(L+\eta)}).
\end{equation}
Moreover \(\|\delta_t\|_\infty=O_\nu(t^{-2})\).
More precisely, the remainder satisfies
\[
\|\rho_{L,t}^{(\eta)}\|_{\infty}
+\|\rho_{L,t}^{(\eta)}\|_{L^1(\omega)}
\le
C_{L,\eta}\,\EE|X_c|^{L+\eta}\,t^{-(L+\eta)},
\]
where \(C_{L,\eta}\) depends only on \(L,\eta\) and the fixed Cayley-mode
envelopes.  For \(0<\eta<1\), this is a genuine fractional-moment
improvement; for \(\eta=1\), it is the usual \((L+1)\)-moment estimate.

Consequently, if \(N\ge2\), \(L=2N-2\), and
\(\EE|X_c|^{2N-2+\eta}<\infty\), then
\begin{equation}\label{eq:positive-margin-retained-order-theorem}
D_{\mathrm{KL}}(P_t*\nu\|P_t(\cdot-\bar{\gamma}))
=\sum_{m=2}^{N}A_{2m}(\nu)t^{-2m}
+O_\nu(t^{-(2N+\eta)}).
\end{equation}
\label{eq:explicit-entropy-constant-bound}
\end{lemma}

\begin{proof}
The quotient identity is
\[
\frac{P_t(\bar{\gamma}+ty-\gamma)}{P_t(ty)}
=R_{X_c/t}(y).
\]
Subtract the Taylor polynomial
\(\sum_{n=0}^{L}u_n(y)(X_c/t)^n\) and split the expectation into
\(|X_c|\le t\) and \(|X_c|>t\).  Both sides of the split are estimated by the
same \(L+\eta\) moment.

On \(|X_c|\le t\), the first line of
\eqref{eq:one-dimensional-two-region-taylor-remainder} gives, uniformly in
\(y\),
\[
\left|
\EE\!\left[
\left(R_{X_c/t}(y)-\sum_{n=0}^{L}u_n(y)(X_c/t)^n\right)
\mathbf1_{\{|X_c|\le t\}}
\right]\right|
\le
C_Lt^{-(L+1)}
\EE\!\left[|X_c|^{L+1}\mathbf1_{\{|X_c|\le t\}}\right].
\]
Because \(0<\eta\le1\) and \(|X_c|\le t\), the last term is bounded by
\[
C_Lt^{-(L+\eta)}
\EE\!\left[
|X_c|^{L+\eta}
\left(\frac{|X_c|}{t}\right)^{1-\eta}
\mathbf1_{\{|X_c|\le t\}}
\right]
\le
C_L\EE|X_c|^{L+\eta}\,t^{-(L+\eta)} .
\]
At the endpoint \(\eta=1\) this is precisely the assumed
\((L+1)\)-moment bound; for \(0<\eta<1\) it uses only the lower
\((L+\eta)\)-moment.

On \(|X_c|>t\), the second line of
\eqref{eq:one-dimensional-two-region-taylor-remainder} gives
\[
\left|
\EE\!\left[
\left(R_{X_c/t}(y)-\sum_{n=0}^{L}u_n(y)(X_c/t)^n\right)
\mathbf1_{\{|X_c|>t\}}
\right]\right|
\le
C_Lt^{-L}\EE\!\left[|X_c|^L\mathbf1_{\{|X_c|>t\}}\right].
\]
Since \(|X_c|>t\) implies
\(|X_c|^L\le t^{-\eta}|X_c|^{L+\eta}\), this is also
\[
O_\nu(t^{-(L+\eta)}).
\]
The same uniform bounds give the \(L^1(\omega)\) estimate because
\(\omega\) is a probability measure.  The \(n=0\) term cancels in
\(\delta_t\), and the \(n=1\) term vanishes after integration against
\(\nu\) because \(\EE X_c=0\).  This proves
\eqref{eq:positive-margin-density-remainder}.  The bound
\(\|\delta_t\|_\infty=O_\nu(t^{-2})\) follows from
Lemma~\ref{lem:exact-moment-density-truncation}, since
\(\EE|X_c|^{L+\eta}<\infty\) and \(L\ge2\).

Now put \(L=2N-2\) and
\[
d_t(y)=\sum_{n=2}^{L}\mu_nu_n(y)t^{-n}.
\]
As in the proof of Theorem~\ref{thm:poisson-kl-even-orders},
\[
D_{\mathrm{KL}}(P_t*\nu\|P_t(\cdot-\bar{\gamma}))
=\int_{\RR}\Phi(\delta_t(y))\,\omega(\dd y),
\qquad
\Phi(s)=(1+s)\log(1+s)-s .
\]
The proof of
Proposition~\ref{prop:finite-quotient-jet-implies-kl-coefficient-jet}, with
the stronger bound
\(\|\rho_{L,t}^{(\eta)}\|_\infty+
\|\rho_{L,t}^{(\eta)}\|_{L^1(\omega)}
=O_\nu(t^{-(L+\eta)})\), gives the same finite monomial coefficients as
Theorem~\ref{thm:poisson-kl-even-orders}; every product containing the
quotient remainder is
\(O_\nu(t^{-(L+\eta+2)})=O_\nu(t^{-(2N+\eta)})\), and the Taylor tail is
\(O_\nu(t^{-2N-2})\).  This proves
\eqref{eq:positive-margin-retained-order-theorem}.
\end{proof}

\begin{lemma}[Lower moments do not identify the next even moment]%
\label{lem:lower-moment-realizability-obstruction}
Let \(L\ge4\) be even.  There exist two centred symmetric compactly
supported probability laws \(\nu_+\) and \(\nu_-\) on \(\RR\) such that
\[
  \int x^j\,\nu_+(\dd x)=\int x^j\,\nu_-(\dd x)
  \qquad(0\le j\le L-1),
\]
but
\[
  \int x^L\,\nu_+(\dd x)\ne \int x^L\,\nu_-(\dd x).
\]
In particular, even within compactly supported centred symmetric laws, the
moment vector through order \(L-1\) does not determine the \(L\)th centred
moment.
\end{lemma}

\begin{proof}
Set \(x_i=i-L/2\), \(0\le i\le L\), and define the signed measure
\[
  \sigma=\sum_{i=0}^{L}(-1)^{L-i}\binom{L}{i}\delta_{x_i}.
\]
This is the \(L\)th finite-difference functional translated by \(L/2\).
Hence, for every polynomial \(p\) of degree \(<L\),
\[
  \int p(x)\,\sigma(\dd x)
  =\sum_{i=0}^{L}(-1)^{L-i}\binom{L}{i}p(i-L/2)=0,
\]
whereas
\[
  \int x^L\,\sigma(\dd x)=L!\ne0 .
\]
Because \(L\) is even, the coefficient at \(x_i\) equals the coefficient at
\(-x_i=x_{L-i}\); thus \(\sigma\) is symmetric and has total mass zero.

Let
\[
  \lambda=\frac1{L+1}\sum_{i=0}^{L}\delta_{x_i}.
\]
Choose \(\varepsilon>0\) so small that all weights of
\(\lambda\pm\varepsilon\sigma\) are nonnegative, and put
\(\nu_\pm=\lambda\pm\varepsilon\sigma\).  Since \(\sigma(\RR)=0\), these are
probability laws.  They are symmetric, hence centred.  The moment identities
up to order \(L-1\) follow from the annihilation of lower-degree
polynomials by \(\sigma\), while their \(L\)th moments differ by
\[
  \int x^L\,\nu_+(\dd x)-\int x^L\,\nu_-(\dd x)
  =2\varepsilon L!\ne0 .
\]
\end{proof}

\begin{proposition}[Top-moment dependence of the universal coefficient]%
\label{prop:poisson-kl-moment-threshold-sharpness}
\label{prop:analytic-first-evaluability-obstruction}
For \(N\ge3\), the universal top coefficient \(A_{2N}^{\rm form}\) contains
the nonzero monomial
\[
\kappa_N\,\mathfrak m_2\mathfrak m_{2N-2},
\qquad
\kappa_N=(-1)^N(N-1)2^{-2N+2}\ne0 .
\]
More precisely, with \(L=2N-2\),
\[
  A_{2N}^{\rm form}
  =\kappa_N\,\mathfrak m_2\mathfrak m_L
  +P_N(\mathfrak m_2,\ldots,\mathfrak m_{L-1}),
\]
where \(P_N\) involves no \(\mathfrak m_L\).  Consequently the displayed
universal polynomial \(A_{2N}\) cannot be evaluated from centred moments of
order at most \(2N-3\) on the full class of laws; two compactly supported
centred symmetric laws can have the same such lower moments but different
values of \(A_{2N}\).

This is the algebraic sharpness mechanism in
Theorem~\ref{thm:sharp-moment-equivalence-hierarchy}.  It also proves the
following stronger non-identifiability statement: if the object to be
recovered is the same polynomial \(A_{2N}\) on the full finite-moment class,
then centred moments through order \(2N-3\) do not contain enough
information.  This does not assert a universal renormalized entropy residual
for arbitrary laws outside the coefficient domain.  At the regularly varying
moment boundary, however, the tail class selects a canonical integrated
slowly varying normalization; this case is closed by
Theorem~\ref{cor:main-body-conditional-missing-top-moment-tail}.  The
unconditional analytic tail threshold is the leading variance theorem,
Theorem~\ref{thm:sharp-leading-kl-variance-threshold}.
\end{proposition}

\begin{proof}
Put \(L=2N-2\).  In the formal coefficient
\eqref{eq:universal-entropy-coefficient}, a term involving
\(\mathfrak m_L\) and contributing to total order \(2N=L+2\) must come from
the quadratic part of
\((1+s)\log(1+s)-s\), with the ordered index pairs \((2,L)\) and \((L,2)\).
All terms with three or more factors have total index at least \(2+2+2\) and,
if one factor is \(L\), total index at least \(L+4> L+2\).  Therefore
\[
  A_{2N}^{\rm form}
  =J(2,L)\,\mathfrak m_2\mathfrak m_L
  +P_N(\mathfrak m_2,\ldots,\mathfrak m_{L-1}).
\]
Writing \(L=2m\), the uniform coefficient formula
\eqref{eq:mode-laurent-coefficients} gives
\(q_{2,\pm1}=q_{2,\pm2}=-1/4\) and
\(q_{2m,\pm k}=(-1)^m2^{-2m}\binom{2m-1}{k-1}\).  Hence only the exponents
\(\pm1\) and \(\pm2\) of \(Q_{2m}\) contribute to \(J(2,2m)\), and
\[
\begin{aligned}
J(2,2m)
&=[z^0]Q_2Q_{2m}  \\
&=-\frac12\,(-1)^m2^{-2m}
\left\{\binom{2m-1}{0}+\binom{2m-1}{1}\right\}  \\
&=(-1)^{m+1}m\,2^{-2m}
=(-1)^N(N-1)2^{-2N+2}.
\end{aligned}
\]
This proves the displayed top-moment coefficient.

It remains to justify the stated lower-moment non-identifiability.  Apply
Lemma~\ref{lem:lower-moment-realizability-obstruction} with this
\(L=2N-2\).  The resulting compactly supported symmetric laws \(\nu_+\) and
\(\nu_-\) have the same centred moments through order \(L-1\), in particular
the same positive second moment, but they have different \(L\)th centred
moments.  Since
\[
  A_{2N}(\nu_+)-A_{2N}(\nu_-)
  =\kappa_N\,\mu_2\{\mu_L(\nu_+)-\mu_L(\nu_-)\}
  \ne0,
\]
no rule depending only on centred moments of orders at most \(L-1=2N-3\)
can evaluate \(A_{2N}\) on the full class of laws.
\end{proof}

\begin{theorem}[Regularly varying missing-moment boundary layer]%
\label{cor:main-body-conditional-missing-top-moment-tail}
Let \(N\ge3\), put \(p=2N-2\), and let \(X\) have a centred probability
law \(\nu\) on \(\RR\).  Assume that, for some measurable eventually
positive slowly varying function \(L\) and constants \(c_+,c_-\ge0\) with
\(C:=c_++c_->0\),
\[
 \frac{\mathbb P(X>x)}{x^{-p}L(x)}\longrightarrow c_+,\qquad
 \frac{\mathbb P(X<-x)}{x^{-p}L(x)}\longrightarrow c_-,
 \qquad x\to\infty.
\]
Fix \(x_0>0\) beyond which \(L\) is positive, and assume that
\[
 \ell_L(t):=\int_{x_0}^{t}\frac{L(s)}s\,\dd s\longrightarrow\infty.
\]
Then \(\EE|X|^q<\infty\) for every \(q<p\), whereas
\(\EE|X|^p=\infty\), and the truncated boundary moment satisfies
\[
 M_p(t):=\EE\!\left[|X|^p\mathbf1_{\{|X|\le t\}}\right]
 \sim pC\ell_L(t).
\]
With \(H(t)=D_{\mathrm{KL}}(P_t*\nu\|P_t)\) and
\(\kappa_N=(-1)^N(N-1)2^{-2N+2}\), one has the exact boundary expansion
\[
H(t)=
\sum_{j=2}^{N-1}A_{2j}(\nu)t^{-2j}
+\kappa_N\mu_2\,t^{-2N}M_p(t)
+o(t^{-2N}M_p(t)).
\]
Equivalently,
\[
 \lim_{t\to\infty}\frac{t^{2N}}{M_p(t)}
 \left\{H(t)-\sum_{j=2}^{N-1}A_{2j}(\nu)t^{-2j}\right\}
 =\kappa_N\mu_2,
\]
and
\[
 \lim_{t\to\infty}\frac{t^{2N}}{\ell_L(t)}
 \left\{H(t)-\sum_{j=2}^{N-1}A_{2j}(\nu)t^{-2j}\right\}
 =p(c_++c_-)\kappa_N\mu_2.
\]
Thus the leading boundary layer depends on the sum of the two tail constants
but not on their balance.  In particular, \(L\equiv1\) gives a \(\log t\)
correction, while \(L(t)=(\log t)^{-1}\) gives a \(\log\log t\)
correction.  No density, symmetry, or monotonicity of \(L\) is assumed.
\end{theorem}

\subsection{Sharp leading threshold}

\begin{lemma}[Second-mode transform of a Cauchy translate]%
\label{lem:second-mode-cauchy-translate-transform}
For \(\varepsilon\in\RR\),
\[
\int_{\RR}u_2(y)\{R_\varepsilon(y)-1\}\,\omega(\dd y)
=\frac{4\varepsilon^2}{(4+\varepsilon^2)^2}.
\]
\end{lemma}

\begin{proof}
Since \(\int_{\RR}u_2\,\dd\omega=0\), it is enough to compute
\(\int u_2R_\varepsilon\,\dd\omega\).  We use the same Cayley/Fourier
normalization as in Theorem~\ref{thm:finite-fourier-mode-integral-algorithm}.
Put \(y=\tan\theta\) and \(z=e^{2\ii\theta}\).  Then
\(\omega(\dd y)=\dd\theta/\pi\), Haar integration on the \(z\)-circle is
constant-term extraction, and
\[
  Q_2(z)=u_2(\tan\theta)
  =-\frac14(z+z^{-1}+z^2+z^{-2}).
\]
For \(\varepsilon=0\), \(R_0\equiv1\) and the integral is zero.  Assume
\(\varepsilon\ne0\).  From
\[
R_\varepsilon(\tan\theta)
=\{1-\varepsilon\sin2\theta+\varepsilon^2\cos^2\theta\}^{-1}
\]
one obtains, after multiplying numerator and denominator by \(z\),
\[
R_\varepsilon(\tan\theta)
=
\frac{z}
{\frac{\varepsilon(\varepsilon+2\ii)}4z^2
+(1+\varepsilon^2/2)z
+\frac{\varepsilon(\varepsilon-2\ii)}4}
=
\frac{\alpha}{z-\alpha}-\frac{\beta}{z-\beta},
\]
where
\[
\alpha=-\frac{\varepsilon}{\varepsilon+2\ii},
\qquad
\beta=\frac{2\ii-\varepsilon}{\varepsilon}.
\]
Here \(|\alpha|<1<|\beta|\), so on the unit circle the Laurent expansion is
\[
R_\varepsilon(\tan\theta)
=\sum_{n=0}^{\infty}\beta^{-n}z^n
+\sum_{n=1}^{\infty}\alpha^nz^{-n}.
\]
Hence
\[
\begin{aligned}
\int_{\RR}u_2(y)R_\varepsilon(y)\,\omega(\dd y)
&=[z^0]\,Q_2(z)R_\varepsilon(\tan\theta)\\
&=-\frac14\left(\alpha+\beta^{-1}+\alpha^2+\beta^{-2}\right)\\
&=-\frac14\left\{
-\frac{\varepsilon}{\varepsilon+2\ii}
-\frac{\varepsilon}{\varepsilon-2\ii}
+\frac{\varepsilon^2}{(\varepsilon+2\ii)^2}
+\frac{\varepsilon^2}{(\varepsilon-2\ii)^2}
\right\}  \\
&=-\frac14\left\{
-\frac{2\varepsilon^2}{4+\varepsilon^2}
+\frac{2\varepsilon^2(\varepsilon^2-4)}
{(4+\varepsilon^2)^2}
\right\}  \\
&=-\frac14\left\{
-\frac{16\varepsilon^2}{(4+\varepsilon^2)^2}
\right\}  \\
&=\frac{4\varepsilon^2}{(4+\varepsilon^2)^2}.
\end{aligned}
\]
The last line is an elementary rational simplification; equivalently, the
bracket equals \(-16\varepsilon^2/(4+\varepsilon^2)^2\).  This proves the
claimed transform.
\end{proof}

\begin{theorem}[Sharp analytic variance threshold for the leading KL coefficient]%
\label{thm:sharp-leading-kl-variance-threshold}
Let \(\nu\) have finite mean \(\bar{\gamma}\), and put
\[
H(t)=D_{\rm KL}(P_t*\nu\|P_t(\cdot-\bar{\gamma})).
\]
With the extended nonnegative convention
\eqref{eq:extended-kl-definition} for \(D_{\rm KL}\),
\[
H(t)\in[0,\infty]\qquad(t>0).
\]
Then
\[
\EE X_c^2<\infty
\quad\Longleftrightarrow\quad
\limsup_{t\to\infty}t^4H(t)<\infty .
\]
When the variance is finite,
\[
H(t)<\infty\qquad(t>0)
\]
and
\[
\lim_{t\to\infty}t^4H(t)=\frac{(\EE X_c^2)^2}{8},
\]
and when the variance is infinite,
\[
\lim_{t\to\infty}t^4H(t)=+\infty .
\]
The last limit is meant in the extended sense; no finite-time finiteness of
\(H(t)\) is asserted in the infinite-variance case, and the conclusion
includes the case \(H(t)=+\infty\) along an unbounded set of times.
No compact-support, exponential-tail, or regular-variation hypothesis is
assumed; finite variance is exactly one of the two equivalent conclusions.
This is the unconditional analytic iff asserted in this paper.
\end{theorem}

\begin{proof}
Finite variance is Theorem~\ref{thm:poisson-kl-even-orders} with \(N=2\).
In that case Lemma~\ref{lem:finite-time-poisson-kl-finite} also gives
\(H(t)<\infty\) for every \(t>0\).  Hence the finite-variance limit is an
ordinary real limit.

For the converse, write
\[
h_t=P_t*\nu,\qquad g_t=P_t(\cdot-\bar{\gamma}),\qquad
\delta_t(y)=\frac{h_t(\bar{\gamma}+ty)}{g_t(\bar{\gamma}+ty)}-1 .
\]
Let \(\varphi_t(x)=u_2((x-\bar{\gamma})/t)\).  Since
\(\|\varphi_t\|_\infty=\|u_2\|_\infty\le1\),
Lemma~\ref{lem:bounded-test-pinsker} gives, in the
extended-valued sense,
\begin{equation}\label{eq:bounded-test-pinsker-second-mode}
2H(t)\ge
\left|
\int_{\RR}\varphi_t(x)\{h_t(x)-g_t(x)\}\dd x
\right|^2
=
\left|
\int_{\RR}u_2(y)\delta_t(y)\,\omega(\dd y)
\right|^2 .
\end{equation}
This use of Pinsker involves only the bounded test \(\varphi_t\); it does not
assume that \(H(t)\) is finite.
